\documentclass[11pt]{article}

\usepackage[T1]{fontenc}
\usepackage[utf8]{inputenc}
\usepackage{lmodern}
\usepackage[margin=1in]{geometry}
\usepackage{microtype}
\usepackage{amsmath,amssymb,mathtools,amsthm}
\usepackage{booktabs,tabularx,array}
\usepackage{enumitem}
\usepackage{xcolor}
\usepackage{hyperref}

\hypersetup{
  colorlinks=true,
  linkcolor=blue!55!black,
  citecolor=blue!55!black,
  urlcolor=blue!55!black
}

\newtheorem{definition}{Definition}
\newtheorem{theorem}{Theorem}
\newtheorem{proposition}{Proposition}
\newtheorem{conjecture}{Conjectural theorem target}
\theoremstyle{remark}
\newtheorem{remark}{Remark}

\newcommand{\Spec}{\operatorname{Spec}}
\newcommand{\Var}{\operatorname{Var}}
\newcommand{\Cov}{\operatorname{Cov}}
\newcommand{\R}{\mathbb{R}}
\newcommand{\cA}{\mathcal{A}}
\newcommand{\cF}{\mathcal{F}}
\newcommand{\cH}{\mathcal{H}}

\title{Addressing the Yang--Mills Existence and Mass-Gap Problem\\
\large Through Relational Rigidity and Phenomenon-Preserving Representations}
\author{}
\date{}

\begin{document}

\maketitle

\begin{abstract}
The phenomenon-preserving interpretation of physical representation suggests a
concrete reformulation of the Yang--Mills existence and mass-gap problem.
Rather than privileging gauge potentials, particles, or a particular regulator,
one seeks a coherent continuum family of gauge-invariant observable relations
and proves that local non-Abelian dynamics generates a strictly positive form
of infrared relational rigidity.  This rigidity must be transported through
gauge-invariant coarse-graining, the continuum and infinite-volume limits, and
Osterwalder--Schrader reconstruction to the spectrum of the physical
Hamiltonian.  The programme developed here isolates three theorem targets: a
rigidity-transfer theorem, an RG-to-rigidity theorem, and a Yang--Mills
infrared-basin-entry theorem.  The first is a conditional spectral principle;
the second can be calibrated in restricted models; the third contains the
genuine four-dimensional frontier.  No result in this essay constitutes a
proof of the Clay problem; the purpose is to identify a noncircular proof
architecture and the exact generative theorem still required.
\end{abstract}

\section{The Clay problem under the relational interpretation}

The official target is to prove, for every compact simple gauge group $G$, that
a nontrivial quantum Yang--Mills theory exists on $\R^4$ and has a physical
Hamiltonian $H$ satisfying
\begin{equation}
  \Spec(H)\subseteq\{0\}\cup[\Delta,\infty),
  \qquad \Delta>0.
  \label{eq:clay-target}
\end{equation}
The construction must satisfy sufficiently strong axiomatic quantum-field
theoretic conditions \cite{jaffe-witten}.  The difficulty includes both the
existence of a nontrivial four-dimensional continuum theory and a gap uniform
in the infinite-volume limit.

Under the relational ontology, the problem becomes:
\begin{quote}
Construct a coherent continuum family of gauge-invariant observable relations,
and prove that local non-Abelian dynamics generates a strictly positive form of
infrared relational rigidity that survives every representation change needed
to reconstruct the physical Hamiltonian.
\end{quote}

This divides the problem into three logically distinct obligations:
\begin{equation}
\boxed{
\begin{aligned}
\text{Generation:}\quad&
  \text{local non-Abelian constraints generate infrared rigidity};\\
\text{Transport:}\quad&
  \text{the rigidity survives changes of scale and representation};\\
\text{Realization:}\quad&
  \text{the coherent limiting structure reconstructs a physical QFT}.
\end{aligned}}
\label{eq:three-obligations}
\end{equation}

\section{A gauge-invariant pre-Hamiltonian rigidity object}

The mass gap is naturally stated in Hamiltonian language, but the Hamiltonian
is available only after the continuum theory has been constructed.  A proof
therefore needs a pre-Hamiltonian manifestation of the same physical rigidity.

Let
\begin{equation}
  \mathfrak R_{a,L}
  =
  \left(
    \cA^{\mathrm{inv}}_{a,L},
    \omega_{a,L},
    \Theta,
    \tau,
    \mathcal L
  \right)
  \label{eq:regulated-package}
\end{equation}
denote a regulated Yang--Mills representation.  Here
\begin{itemize}[leftmargin=2em]
  \item $\cA^{\mathrm{inv}}_{a,L}$ is a gauge-invariant observable algebra;
  \item $\omega_{a,L}$ is the Wilson-lattice state;
  \item $\Theta$ is Euclidean time reflection;
  \item $\tau$ is Euclidean translation;
  \item $\mathcal L$ records cutoff, volume, and renormalization maps.
\end{itemize}

For a centered gauge-invariant observable $A$, define the Euclidean temporal
correlation
\begin{equation}
  C^A_{a,L}(t)
  =
  \omega_{a,L}\!\left(A^*\tau_t(A)\right),
  \qquad
  \omega_{a,L}(A)=0.
  \label{eq:regulated-correlation}
\end{equation}
Reflection positivity gives this expression a spectral interpretation after
Osterwalder--Schrader reconstruction.

\begin{definition}[Uniform physical relational rigidity]
A family $\{\mathfrak R_{a,L}\}$ has uniform physical relational rigidity with
rate $m_*>0$ if, for every $m<m_*$ and every $A$ in a
convergence-determining gauge-invariant class,
\begin{equation}
  \left|C^A_{a,L}(t)\right|
  \leq K_A e^{-mt},
  \label{eq:uniform-rigidity}
\end{equation}
where $t$ is measured in physical units, $K_A$ is stable under the required
renormalization of $A$, the estimate is uniform in $L$, it persists along the
continuum trajectory $a\to0$, and the observable class becomes cyclic or
spectrally complete after reconstruction.
\end{definition}

This property is formulated without treating $A_\mu$, a particular gauge
fixing, perturbative particles, or a finite-lattice Hamiltonian as fundamental.
It records the persistence of physical distinctions in a gauge-invariant
observable language.

\section{A conditional rigidity-transfer theorem}

The first exact theorem target is a transport result from regulated Euclidean
rigidity to the continuum physical spectrum.

\begin{theorem}[Relational Rigidity Transfer; conditional form]
Suppose that:
\begin{enumerate}[label=(\roman*),leftmargin=2.3em]
  \item regulated gauge-invariant Schwinger functions converge to a nontrivial
        continuum family;
  \item the limiting family satisfies the Osterwalder--Schrader axioms;
  \item uniform physical relational rigidity holds with rate $m_*>0$;
  \item the observables satisfying the estimate generate a dense subspace of
        the reconstructed physical Hilbert space orthogonal to the vacuum.
\end{enumerate}
Then the reconstructed Hamiltonian obeys
\begin{equation}
  \Spec(H)\subseteq\{0\}\cup[m_*,\infty).
  \label{eq:transferred-gap}
\end{equation}
\end{theorem}

\begin{proof}[Proof mechanism]
Osterwalder--Schrader reconstruction identifies the limiting Euclidean
correlation with
\begin{equation}
  C^A(t)
  =
  \langle\psi_A,e^{-tH}\psi_A\rangle,
  \qquad \psi_A=A\Omega .
  \label{eq:os-semigroup}
\end{equation}
The spectral theorem gives
\begin{equation}
  C^A(t)
  =
  \int_{[0,\infty)}e^{-tE}\,d\mu_A(E).
  \label{eq:spectral-measure}
\end{equation}
If $\mu_A$ assigned positive weight to an interval below $m_*$, its contribution
would decay more slowly than $e^{-mt}$ for some $m<m_*$, contradicting
\eqref{eq:uniform-rigidity}.  Hence every centered vector generated by the
designated observables has spectral support in $[m_*,\infty)$.  Density excludes
an additional physical spectral subspace in $(0,m_*)$.
\end{proof}

\begin{remark}
This theorem is conditional and belongs to the established family of
spectral--clustering principles.  Its role here is architectural: it states
exactly which rigidity information must be transported through the continuum
construction.  It does not generate the required estimate.
\end{remark}

\section{The real problem: generation of the rigidity estimate}

The difficult step is now isolated:
\begin{quote}
Derive a volume-uniform exponential-decay estimate, in physical units, from
local Yang--Mills dynamics along the four-dimensional continuum trajectory.
\end{quote}

At fixed strong lattice coupling, cluster expansions, stochastic dynamics, and
functional inequalities can yield infinite-volume limits and exponential decay.
Recent Poincar\'e and log-Sobolev approaches give such results for lattice
Yang--Mills in specified strong-coupling regimes \cite{shen-zhu-zhu,nissim}.
These are genuine lattice correlation-decay results.  They do not reach the
weak bare-coupling limit $\beta\to\infty$ associated with the four-dimensional
continuum construction.

The missing inference is therefore not the formal implication
\begin{equation}
  \text{finite lattice gap}
  \longrightarrow
  \text{continuum gap}.
\end{equation}
It is the dynamical statement
\begin{equation}
\boxed{
  \text{asymptotically free ultraviolet theory}
  \longrightarrow
  \text{controlled entry into an infrared rigidity regime}.}
  \label{eq:missing-inference}
\end{equation}

\section{A two-regime proof architecture}

The representation ontology suggests separating ultraviolet transport from
infrared rigidity generation.

\subsection{Regime I: ultraviolet coherence transport}

Beginning with a fine Wilson lattice, one seeks to
\begin{enumerate}[leftmargin=2em]
  \item perform gauge-covariant block transformations;
  \item integrate out short-distance degrees of freedom;
  \item transport a convergence-determining gauge-invariant observable algebra;
  \item preserve reflection positivity or prove an adequate recovery theorem;
  \item control small-field and large-field sectors separately;
  \item retain explicit dependence on lattice spacing and running coupling.
\end{enumerate}

The intended chain is
\begin{equation}
  \mathfrak R_a
  \longrightarrow
  \mathfrak R_{ba}
  \longrightarrow
  \mathfrak R_{b^2a}
  \longrightarrow\cdots\longrightarrow
  \mathfrak R_{\ell_*}.
  \label{eq:rg-chain}
\end{equation}
Balaban's four-dimensional analysis provides rigorous small-field
renormalization-group structure, but it does not complete the continuum
construction or establish the infrared gap \cite{balaban}.

\subsection{Regime II: infrared rigidity generation}

At a dynamically generated physical scale $\ell_*$, the effective
gauge-invariant representation should enter a domain in which one can prove
conditional mixing, a block functional inequality, transfer-operator
contraction, effective coercivity, or another local positive inequality that
implies exponential correlation decay.

A candidate block inequality is
\begin{equation}
  \Var_{\mu_*}(F)
  \leq
  \frac{1}{\lambda_*}
  \sum_B
  \mathbb E_{\mu_*}
  \left[
    \Var\!\left(F\mid\cF_{B^c}\right)
  \right],
  \qquad \lambda_*>0,
  \label{eq:block-poincare}
\end{equation}
together with sufficiently weak inter-block influence.  The desired consequence
is
\begin{equation}
  \left|\Cov_{\mu_*}(A,B)\right|
  \leq
  C_{A,B}
  e^{-c\,d(A,B)/\ell_*},
  \label{eq:block-clustering}
\end{equation}
giving the physical rate
\begin{equation}
  m_* = \frac{c}{\ell_*}.
  \label{eq:physical-rate}
\end{equation}

A Poincar\'e inequality for an auxiliary Langevin dynamics is not itself the
quantum Hamiltonian mass gap.  It is an analytic engine for deriving physical
correlation decay.  The rigidity-transfer theorem must then carry that decay to
the OS-reconstructed Hamiltonian.

\section{The central missing theorem: infrared basin entry}

The interpretation identifies a sharper frontier target than the undivided
instruction to ``prove the mass gap.''

\begin{conjecture}[Yang--Mills Infrared Basin Entry]
The gauge-covariant renormalization flow of four-dimensional pure Yang--Mills
theory, starting on the asymptotically free continuum trajectory, reaches at a
finite physical scale $\ell_*$ a gauge-invariant effective representation
satisfying a verifiable local rigidity criterion:
\begin{equation}
  \mathfrak R_{a,g(a)}
  \xrightarrow[\ a\to0\ ]{\text{controlled RG}}
  \mathfrak R_{\ell_*}
  \in
  \mathfrak B_{\mathrm{rigid}}.
  \label{eq:basin-entry}
\end{equation}
Here $\mathfrak B_{\mathrm{rigid}}$ is defined by explicit local inequalities
and does not assume a mass gap or exponential clustering.
\end{conjecture}

A proof would have to show:
\begin{enumerate}[leftmargin=2em]
  \item $\ell_*$ remains finite and nonzero in physical units;
  \item the effective representation retains the correct Yang--Mills observable
        content;
  \item reflection positivity or a sufficient substitute survives;
  \item the rigidity criterion does not contain clustering among its premises;
  \item large-field and topological sectors are controlled;
  \item the constants are uniform in the original volume;
  \item pulling the bound back through the RG chain does not make it vanish.
\end{enumerate}

Perturbative renormalization suggests a scale of the schematic form
\begin{equation}
  \ell_*
  \sim
  \mu^{-1}
  \exp\!\left(\frac{1}{2b_0g^2(\mu)}\right),
  \label{eq:transmutation-scale}
\end{equation}
but the running-coupling formula does not prove that the full theory enters a
massive infrared basin at that scale.  The basin-entry theorem is where
dimensional transmutation must become a nonperturbative mathematical statement.

\section{The complete representation atlas}

The proposed proof chain is
\begin{equation}
\begin{aligned}
  \mathfrak R_{\mathrm{Wilson}}
  &\xrightarrow{\text{gauge-invariant RG}}
  \mathfrak R_{\mathrm{block}}
  \xrightarrow{\text{basin entry}}
  \mathfrak R_{\mathrm{rigid}}\\
  &\xrightarrow{\text{uniform clustering}}
  \mathfrak R_{\mathrm{Euclidean}}
  \xrightarrow{\text{continuum coherence}}
  \mathfrak R_{\mathrm{OS}}
  \xrightarrow{\text{reconstruction}}
  \mathfrak R_{\mathrm{Hamiltonian}}.
\end{aligned}
\label{eq:representation-atlas}
\end{equation}

Each arrow requires an independent certificate.

\begin{table}[ht]
\centering
\small
\renewcommand{\arraystretch}{1.2}
\begin{tabularx}{\textwidth}{>{\raggedright\arraybackslash}p{0.34\textwidth}X}
\toprule
\textbf{Representation change} & \textbf{Required certificate} \\
\midrule
Wilson $\to$ block variables
  & Preservation of gauge-invariant observable content \\
Block variables $\to$ rigid regime
  & Noncircular local coercivity or mixing theorem \\
Rigid regime $\to$ clustering
  & Quantitative exponential decay in physical units \\
Regulated theory $\to$ continuum
  & Tightness, renormalization, nontriviality, and symmetry recovery \\
Euclidean limit $\to$ OS theory
  & Reflection positivity and the remaining OS axioms \\
OS theory $\to$ Hamiltonian gap
  & Spectral completeness and rigidity transfer \\
\bottomrule
\end{tabularx}
\caption{The certificates required along the proposed proof chain.}
\label{tab:certificates}
\end{table}

The ontology prevents a property established in one chart from being silently
treated as though it had already survived every subsequent transformation.

\section{The role of finite Wilson/OS constructions}

A finite Wilson/Osterwalder--Schrader $SU(2)$ construction can serve as a
calibrated representation package:
\begin{equation}
  \mathfrak R^{SU(2)}_{a,L}
  \longrightarrow
  \cH_{a,L}
  \longrightarrow
  H_{a,L}.
  \label{eq:finite-package}
\end{equation}
The Clay-relevant task is to expose the complete dependence
\begin{equation}
  \Delta_{a,L}
  =
  \Delta
  \left(
    a,L,\beta,	ext{geometry},	ext{observable sector}
  \right)
  \label{eq:gap-dependence}
\end{equation}
and determine whether
\begin{equation}
  \liminf_{\substack{a\to0\\L\to\infty}}
  \frac{\Delta^{\mathrm{lattice}}_{a,L}}{a}
  >0.
  \label{eq:physical-gap-limit}
\end{equation}

The finite construction becomes relevant to the continuum problem if it
provides at least one of the following:
\begin{itemize}[leftmargin=2em]
  \item a physical-unit lower bound uniform in volume;
  \item a block-local inequality stable under renormalization;
  \item an exact comparison between transfer spectra and gauge-invariant
        correlation decay;
  \item a geometric coercivity estimate whose constants survive continuum
        scaling;
  \item control of large-field corrections to a heat-kernel or small-field
        representation.
\end{itemize}
Its role is therefore not merely to display a finite eigenvalue separation, but
to serve as a laboratory for discovering a transportable rigidity certificate.

\section{What may change and what must survive}

The proof may change gauge coordinates, lattice variables, gauge fixing,
regulators, block variables, local parametrizations, particle interpretations,
and apparent microscopic ontology.  It must preserve or recover
\begin{itemize}[leftmargin=2em]
  \item gauge-invariant observable content;
  \item reflection positivity and physical positivity;
  \item locality or causal reconstructibility;
  \item four-dimensional continuum covariance;
  \item asymptotic-freedom scaling;
  \item nontriviality of the limiting theory;
  \item completeness of the physical sectors;
  \item a positive rigidity rate in physical units.
\end{itemize}
This permits a search for nonclassical variables without changing the target
theory, provided every representational bridge is certified.

\section{The irreducible core and theorem sequence}

The ontology does not dissolve the Clay problem.  It exposes its irreducible
core:
\begin{equation}
\boxed{
\begin{gathered}
\text{Prove that four-dimensional non-Abelian RG flow generates}\
\text{a finite physical scale at which gauge-invariant relational}\
\text{influence becomes uniformly contractive, and prove that this}\
\text{contractivity survives continuum reconstruction.}
\end{gathered}}
\label{eq:irreducible-core}
\end{equation}

The research programme should pursue three theorem targets in order:
\begin{enumerate}[leftmargin=2em]
  \item \textbf{Rigidity Transfer Theorem:} uniform decay on a complete
        regulated observable class implies a continuum Hamiltonian gap.
  \item \textbf{RG-to-Rigidity Theorem:} entry into an explicitly defined
        functional-inequality or mixing domain implies uniform physical
        relational rigidity.
  \item \textbf{Yang--Mills Basin Entry Theorem:} the four-dimensional pure
        Yang--Mills RG trajectory enters that domain at a finite dynamically
        generated scale.
\end{enumerate}

The first target is tractable as a precise transport theorem.  The second can be
proved and stress-tested in restricted models.  The third is the genuine
frontier theorem.  Its validity depends on formulating $\mathfrak
B_{\mathrm{rigid}}$ through independently verifiable local inequalities rather
than by hiding the desired mass gap in its definition.

\begin{thebibliography}{99}

\bibitem{jaffe-witten}
A.~Jaffe and E.~Witten,
\emph{Quantum Yang--Mills Theory},
in \emph{The Millennium Prize Problems},
Clay Mathematics Institute and American Mathematical Society, 2006.
\url{https://www.claymath.org/wp-content/uploads/2022/06/yangmills.pdf}.

\bibitem{balaban}
T.~Ba\l aban,
\emph{Renormalization Group Approach to Lattice Gauge Field Theories I:
Generation of Effective Actions in a Small Field Approximation and a Coupling
Constant Renormalization in 4D},
Communications in Mathematical Physics \textbf{109} (1987), 249--301.

\bibitem{shen-zhu-zhu}
H.~Shen, R.~Zhu, and X.~Zhu,
\emph{Langevin Dynamics of Lattice Yang--Mills and Applications},
arXiv:2204.12737, 2022.
\url{https://arxiv.org/abs/2204.12737}.

\bibitem{nissim}
R.~Nissim,
\emph{$\mathrm{U}(N)$ Lattice Yang--Mills in the 't Hooft Regime},
arXiv:2510.22788, 2025.
\url{https://arxiv.org/abs/2510.22788}.

\bibitem{chatterjee}
S.~Chatterjee,
\emph{Yang--Mills for Probabilists},
arXiv:1803.01950, 2018.
\url{https://arxiv.org/abs/1803.01950}.

\bibitem{osterwalder-schrader}
K.~Osterwalder and R.~Schrader,
\emph{Axioms for Euclidean Green's Functions},
Communications in Mathematical Physics \textbf{31} (1973), 83--112;
\textbf{42} (1975), 281--305.

\end{thebibliography}

\end{document}
